\section{Efficient VLAs and Action Generation}
\label{sec:study}

\paragraph{Where efficient VLAs save cost.}
\cref{tab:study} groups the efficiency work in the literature matrix (18 sources, each read from the paper itself; delivered as a spreadsheet) by the part of the inference loop it changes.
Three assumptions recur.
Token and layer reduction assumes redundancy: consecutive frames share most visual content~\cite{xu2025vlacache}, or easy situations need fewer layers~\cite{yue2024deer}; both save compute only in the language model, not in the vision encoder.
Chunked and asynchronous execution assumes that actions predicted from an older observation remain valid; SmolVLA's asynchronous mode shortened task completion by 30\% but lowered success from 78.3\% to 73.3\%~\cite{shukor2025smolvla}, and A2C2 adds a trained correction head to restore reactivity~\cite{sendai2025a2c2}.
Systems work assumes that the model stays numerically the same, which an export can violate~\cite{islam2026onnx}.
Two open questions follow for this study.
First, speed-ups are usually reported next to aggregate success rates, yet a non-significant difference is not evidence of equivalence~\cite{islam2026onnx}; this study therefore fixes a non-inferiority margin in advance and tests it on paired episodes.
Second, most latencies are measured on GPUs with native bfloat16 support (RTX 4090, A100, A40), so it is unclear which bottleneck dominates on older or embedded GPUs.

\begin{table}[t]
\centering\scriptsize
\setlength{\tabcolsep}{2.5pt}
\begin{tabular}{@{}p{1.55cm}p{2.05cm}p{3.95cm}@{}}
\toprule
Lever & Examples & Reported effect (hardware) \\
\midrule
Smaller model & SmolVLA~\cite{shukor2025smolvla}, TinyVLA~\cite{wen2025tinyvla} & 0.45B params, LIBERO avg 87.3; 14 vs 292\,ms per action against OpenVLA (A6000) \\
Fewer tokens or layers & VLA-Cache~\cite{xu2025vlacache}, EfficientVLA~\cite{yang2025efficientvla}, DeeR~\cite{yue2024deer} & $1.63\times$ (OpenVLA, LIBERO, 4090); $1.93\times$ at 28.9\% FLOPs, $-0.6$ SR (A40); LLM FLOPs $5.2$--$6.5\times$ lower (trained) \\
Action decoding & FAST~\cite{pertsch2025fast}, OFT~\cite{kim2025oft}, CogACT~\cite{li2024cogact} & up to $5\times$ faster training but 750 vs 100\,ms per chunk; $26\times$ throughput; VLM runs once, small DiT denoises \\
Execution & RTC~\cite{black2025rtc}, A2C2~\cite{sendai2025a2c2} & tolerates $>$300\,ms delay at +21\,ms per chunk; 4.7\,ms head per step \\
Systems & Realtime-VLA~\cite{ma2025realtime}, ONNX~\cite{islam2026onnx}, Coda~\cite{tang2026coda} & $\pi_0$ 106.5$\to$27.3\,ms (4090); 40\% lower median latency with an SR drop; 2-step flow with a learned corrector \\
\bottomrule
\end{tabular}
\caption{Efficiency levers in the literature matrix (selection). Numbers as reported by each paper.}
\label{tab:study}
\end{table}

\paragraph{Diffusion versus flow matching.}
Both families learn to transform Gaussian noise $\boldsymbol\epsilon$ into an action chunk $\mathbf A$ conditioned on the observation $\mathbf o$.
Diffusion heads (Diffusion Policy~\cite{chi2023diffusion}, CogACT~\cite{li2024cogact}) use the DDPM forward process $\mathbf A_k=\sqrt{\bar\alpha_k}\,\mathbf A+\sqrt{1-\bar\alpha_k}\,\boldsymbol\epsilon$ and train a network to predict the injected noise~\cite{ho2020ddpm},
\begin{equation}
\mathcal L_\text{DDPM}=\mathbb E_{k,\boldsymbol\epsilon}\big\|\boldsymbol\epsilon-\boldsymbol\epsilon_\theta(\mathbf A_k,k\mid\mathbf o)\big\|^2 .
\end{equation}
Flow-matching heads ($\pi_0$~\cite{black2024pi0}, SmolVLA) interpolate linearly between noise and data and regress the constant velocity of that path~\cite{lipman2023flow} (\cref{eq:fm}).
The two differ in four respects.
(i) \emph{Target}: noise versus velocity.
(ii) \emph{Path}: the linear interpolation gives each training pair a constant velocity, which yields straighter sampling trajectories than diffusion paths and favours coarse ODE steps~\cite{lipman2023flow}; on RoboTwin, Coda reports $\pi_{0.5}$ at 71.3\% with 3 flow steps against 71.5\% with 10, and 59.8\% with one~\cite{tang2026coda}.
(iii) \emph{Sampler}: Diffusion Policy samples with 100 DDPM steps in simulation and with 10 DDIM~\cite{song2021ddim} steps on the real robot, CogACT with 10 DDIM steps and classifier-free guidance, while flow models integrate the ODE with 10 forward-Euler steps.
(iv) \emph{Architecture}: Diffusion Policy repeats its whole denoising network (9--80M parameters) on visual features computed once; CogACT, $\pi_0$ and SmolVLA run a large VLM once and repeat only a head (a DiT conditioned on one cognition feature, or an action expert that attends to the cached VLM keys and values).
In the last three, the cost of one step is set by the size of the head and, at batch size one, by kernel-launch overhead, which \cref{sec:analysis} measures.
